What is the Eisenberg Noe Model for financial networks

Financial institutions \( N = \{ 1, 2, \ldots, n \} \)
Society is additional node \( 0 \)
Nominal liability matrix: \( (\bar{p}_{ij})_{i,j=0,1,2,\ldots,n} \)
\(
\bar{p}_i = \sum_{j=0}^{n} \bar{p}_{ij}
\)

Relative liabilities:
\[
a_{ij} = 
\begin{cases} 
\frac{\bar{p}_{ij}}{\bar{p}_i}, & \bar{p}_i > 0, \\ 
0, & \bar{p}_i = 0.
\end{cases}
\]

Fixed point / equilibrium:
Liability payments
=
minimum of promised payments and
available resources

INCLUDE IMAGE FROM 472

Call the promised payments of entity \( i \), \( \bar{p}_i \), and \( p_j(x) \) the amount of capital \( i \) received by \( j \)
then \( a_{ji} \) is the fractional amount of wealth that should be paid in case that not all promises are kept.
Additionally \( x_i \) represents the initial capital of \( i \).
\(p_i(x) = \bar{p}_i \wedge \left( \sum_{j=0}^n p_j(x) a_{ji} + x_i \right), \quad i = 1, 2, \ldots, n \)
\( p_0(x) = \bar{p}_0 \)

Denote the own founds of \( i \) by \( e_i \) then
\( e_i(x) = \sum_{j \neq i} p_j(x) a_{ji} + x_i - \bar{p}_i, \quad i = 0, 1, 2, \ldots, n \)

Let \( X \in L^0(\mathbb{R}^n) \) be the random endowment of the financial institutions at time \( t = 1 \) before market clearing.
If we focus on the wealth of the society, then the relevant stochastic outcome is provided by the random field
\( Y_k := e_0(X + k), \quad k \in \mathbb{R}^n. \)


What is a set-valued systemic risk measure? State elementary properties. How are
cash-invariant allocation rules related to regulatory prices of capital for different
entities?

Systemic risk is measured by the set of allocations of additional capital that lead to acceptable outcomes.

Letting \(\mathcal{P}(\mathbb{R}^n; \mathbb{R}_+) := \{ B \subseteq \mathbb{R}^n \mid B = B + \mathbb{R}^n_+ \}\) 
be the collection of upper sets with ordering cone \(\mathbb{R}_+\), we call the function 

\(R: \mathcal{Y} \times \mathbb{R}^n \to \mathcal{P}(\mathbb{R}^n; \mathbb{R}_+^n)\)
a systemic risk measure, if for some acceptance set \(\mathcal{A} \subseteq \mathcal{X}\) of a scalar monetary risk measure:

\[ R(Y; k) = \{ m \in \mathbb{R}^n \mid Y_{k+m} \in \mathcal{A} \}. \]

Elementary Properties:
1. Cash-invariance:  \( R(Y; k) + m = R(Y; k - m) \)
2. Monotonicity:  \( Y \geq Z \implies (\forall k \in \mathbb{R}^n: R(Y; k) \supseteq R(Z; k)) \)


What are factor models? Provide examples.

The value \( V_t \) of a financial position at time \( t \) is explained by the values 
\( Z_t^1, Z_t^2, \ldots, Z_t^m \) of \( m \) risk factors at time \( t \):

\[
V_t = v(Z_t^1, Z_t^2, \ldots, Z_t^m)
\]

for a deterministic function \( v = v(z^1, \ldots, z^m) \).

Factor models are a common approach in practice:
- They are often a consequence of pricing formulas in market models.
- Or, they are approximations of market values in multivariate regression models.

They can be used for simplified calculations based on the distributions of risk factors.

Examples are 

Value of a call option in the Black-Scholes model:

Option with maturity \( T \) and strike \( K \)
Stock with price process \( (S_t)_{t \geq 0} \)
\(
C_t = v(S_t, T - t) \quad \text{with} \quad v(x, \tau) = x \Phi(d_+(x, \tau)) - e^{-r \tau} K \Phi(d_-(x, \tau))
\)

Value of a fixed-rate coupon bond:
Coupon payments \( c_1, \ldots, c_m \) at dates \( T_1, \ldots, T_m \), \( N \) face value
\( P(t, T_i) \) price of a zero bond with face value 1 and maturity \( T_i \) at time \( t \)
\(
V_t = v(P(t, T_1), \ldots, P(t, T_m)) \quad \text{for} \quad v(z^1, \ldots, z^m) = N \cdot z^m + \sum_{i: T_i \geq t} c_i \cdot z^i
\)

Present value of a cashflow stream:
Compounded Spot Rates \( r_t^c(s), s \ge t \)
Deterministic payments \( A_1, \ldots, A_m \) at dates \( T_1, \ldots, T_m \)
\(
V_t = v(r_t^c(T_1), \ldots, r_t^c(T_m)) \quad \text{for} \quad v(z^1, \ldots, z^m) = \sum_{i: T_i \geq t} A_i (1 + z^i)^{-(T_i - t)}
\)

Explain the Delta-normal approximation.
AV@R?

Assumption:
\((\Delta Z^1, \Delta Z^2)^\top \sim \mathcal{N}_2(h\mu, h\Sigma)\)

Conclusions:
\(\mathbb{E}[\Delta Z^i] = h\mu_i, \, \text{Var}[\Delta Z^i] = h\sigma_i^2, \, i = 1, 2, \, \text{Cov}[\Delta Z^1, \Delta Z^2] = h\sigma_{1,2}\)
\(\Delta V\) is approximately normally distributed:
\(\mathbb{E}[\Delta V] \approx h(d_1\mu_1 + d_2\mu_2), \, \text{Var}[\Delta V] \approx h(d_1^2\sigma_1^2 + d_2^2\sigma_2^2 + 2d_1d_2\sigma_{1,2}).\)

Value at Risk of the profit and loss \(\Delta V\):
\(
\text{V@R}_\lambda(\Delta V) = \mathbb{E}[-\Delta V] - \Phi^{-1}(\lambda)\sqrt{\text{Var}(\Delta V)}
\)
Here \(\Phi\) denotes the cumulative distribution function of the standard normal distribution.

Similar formula for Average Value at Risk:
\(
\text{AV@R}_\lambda(\Delta V) = \mathbb{E}[-\Delta V] + \frac{1}{\lambda}\varphi(\Phi^{-1}(\lambda))\sqrt{\text{Var}(\Delta V)}
\)
Here \(\varphi\) denotes the density of the standard normal distribution.

This method is called Delta-Normal approximation:


Explain that distribution-based risk measures define both functionals on measur-
able functions (or elements of the equivalence classes in L0 ) and on probability
measures.

Recall that a risk measure \(\rho : \mathcal{X} \to \mathbb{R}\) is called distribution-based, if
\(
\rho(X) = \rho(Y),
\)
whenever \(\mathbb{P}^X = \mathbb{P}^Y\) with respect to a reference measure \(\mathbb{P}\).
In this case, the risk measure defines a map
\(
\mathcal{R}_{\rho}(\mu) := \rho(X) \text{ with } \mathbb{P}^X = \mu
\)
on Borel probability measures \(\mathcal{M}\).
This motivates to consider risk measures at the level of distributions as statistical functionals 
\(\mathcal{R}_{\rho} : \mathcal{M} \to \mathbb{R}\).


Motivate in the iid-setting that backtesting might be formulated in terms of em-
pirical scores. This motivates a focus on elicitable functionals.


Background
Validation of quantitative risk models requires backtesting in order to assess the accuracy of the forecasting methodology.
Backtesting is the activity of comparing periodically the point predictions of a risk measure with the realized value of the variable under interest.

A canonical implementation of backtesting
Consider a statistical functional \( T = \mathcal{R}_{\rho} \), and let \( x_1, \ldots, x_n \) be a time series of point forecasts of \( T \).
Let \( y_1, \ldots, y_n \) denote the corresponding time series of the variable of interest \( Y \) with distribution \( \mu \in \mathcal{M} \).
The performance is assessed 'good' if - for a given scoring function \( S \) - the average score is 'small':

\(
\hat{s} = \frac{1}{n} \sum_{i=1}^{n} S(x_i, y_i)
\)

Note that for fixed \( x \) and large \( n \), the average score equals approximately the expectation

\(
\int S(x, y) \mu(dy).
\)

This criterium is only meaningful if the statistical functional \( T = \mathcal{R}_{\rho} \) is elicitable, i.e., 
defined in a way that is consistent with the scoring function.


Define the notion of elicitable functionals.

A scoring function \( S : \mathbb{R}^2 \to [0, \infty) \) satisfies:

\( S(x, y) \geq 0 \), and \( S(x, y) = 0 \iff x = y \),
\( x \mapsto S(x, y) \) increasing for \( x > y \) and decreasing for \( x < y \),
\( x \mapsto S(x, y) \) continuous.

Elicitable Statistical Functional
A statistical functional \( T : \mathcal{M} \to \mathbb{R} \) is elicitable on \( \mathcal{M} \), 
if there exists a scoring function \( S \) such that for each \( \mu \in \mathcal{M} \):

\(
\int S(x, y) \mu(dy) < \infty,
\)

\(
T(\mu) = \arg \min_{x \in \mathbb{R}} \int S(x, y) \mu(dy).
\)

(Elicitability)
The expectation \( T(\mu) = \int y \mu(dy) \) is elicitable with respect to the scoring function \( S(x, y) = (y - x)^2 \).


What is the \( ψ \)-weak topology?
Let \(\psi : \mathbb{R} \to [1, \infty)\) be a continuous function which serves as gauge function, and define
\(
\mathcal{M}_1^\psi := \left\{ \mu \in \mathcal{M}_1 : \int \psi \, d\mu < \infty \right\}.
\)

Denote by \(C^\psi(\mathbb{R})\) the space of all continuous functions \(f\) such that \(|f| \leq c|\psi|\).
The \(\psi\)-weak topology on a subset \(\mathcal{M} \subset \mathcal{M}_1^\psi\) is the initial topology generated by

\(
\mathcal{M} \ni \mu \mapsto \int f(y) \mu(dy) \text{ with } f \in C^\psi.
\)
While the weak topology is insensitive to the tails, for unbounded \(\psi\) the \(\psi\)-weak topology is finer 
and does not neglect deviations of measures in the tails.


Provide characterization results for distribution-based risk measures under conditions
such as convex acceptance sets and elicitable functionals

Explain heuristically that the latter characterization b) 
follows (modulo technical details) essentially from a) and Osband’s
observation of convex level sets.

Let \(\rho\) a distribution-based risk measure.
The acceptance set \(\mathcal{A}\) on the level of distributions is defined as:

\(
\mathcal{A} := \{\mu : \mathcal{R}_{\rho}(\mu) \leq 0\}.
\)
The rejection set \(\mathcal{A}^c\) on the level of distributions is given by:
\(
\mathcal{A}^c = \{\mu : \mathcal{R}_{\rho}(\mu) > 0\}.
\)

Observation:
- If \(\rho\) is elicitable, then \(\mathcal{A}\) and \(\mathcal{A}^c\) are convex.
- This follows due to convex level sets on the level of measures, see Osband (1985).

Explain in more detail the methodology of conditional risk measurement in time.

Distortion risk measures are an important class of risk measures that includes VaR, AVaR, RVaR, Wang premium principles, and many other examples as special cases.
Distortion risk measures are a subclass of comonotonic monetary risk measures

We denote in this section by \(X\) a random variable which models losses as its positive part and wealth or gains as its negative part. 
(This convention differs from other sections!)

1. A non-decreasing function \(g : [0, 1] \to [0, 1]\) with \(g(0) = 0\) and \(g(1) = 1\) is called a distortion function
2. On a given probability space \((\Omega, \mathcal{F}, \mathbb{P})\) we define for a distortion function \(g\) a capacity \(c := g \circ \mathbb{P}\)
3. Provided it exists, the Choquet integral
\(
\rho_g(X) := \int X \, dc := \int_{-\infty}^0 (c(X > x) - 1) \, dx + \int_0^{\infty} c(X > x) \, dx
\)
for measurable \(X : \Omega \to \mathbb{R}\) is called the distortion risk measure of \(X\) with distortion function \(g\).


Describe the setup for multinomial backtests.


True loss process:
Stochastic process \(L = (L_t)_{t=1, \ldots, n}\), whose law is unknown
Model process:

Stochastic process \(M = (M_t)_{t=1, \ldots, n}\) with known law
Conditional distribution function \(F_{M_t \mid M_{t-1}, \ldots, M_1}\)

Conditional risk measurement at \(t-1\)
Risk measurements are computed from the model conditionally on \(M_{t-1} = L_{t-1}, \ldots, M_1 = L_1\), i.e., 
the risk measure \(\rho_g\) is applied to the conditional distribution \(F_{M_t \mid L_{t-1}, \ldots, L_1}\)
The result is a function of \(L_{t-1}, \ldots, L_1\) denoted by \(\rho_g^{t-1}(M_t)\)


How can in this situation an AV@R-backtest be reduced to a joint backtest of
multiple V@Rs?

This also includes the question (TODO check if the first question is really this topic):
How does a backtest of multiple V@R lead to a mutinomial test? Describe the
following random variables: number of breached levels at individual time points,
cell counts (by aggregating over time), and the resulting test statistic.

A multinomial backtest for AV@R was developed by Kratz, Lok & McNeil (2018)
Taking an equidistant partition \(\alpha_j := \frac{j}{m+1}\alpha, j = 0, 1, \ldots, m+1\), 
of the interval \([0, \alpha]\), approximate the conditional AV@R as follows:

\(
AV@R_{\alpha}^{t-1}(M_t) = \frac{1}{\alpha} \int_0^{\alpha} V@R_{\lambda}^{t-1}(M_t) d\lambda
\approx \frac{1}{m+1} \left( V@R_{\alpha_1}^{t-1}(M_t) + \cdots + V@R_{\alpha_{m+1}}^{t-1}(M_t) \right)
\)

The backtest considers the random number of breached levels
\(
X_t := \sum_{j=1}^{m+1} \mathbb{1} \left\{ L_t > V@R_{\alpha_j}^{t-1}(M_t) \right\}
\)
and cell counts
\(
O_i = \sum_{t=1}^{n} \mathbb{1} \{ X_t = i \}, \quad i = 0, 1, \ldots, m+1
\)

Compare the methods of Pearon and Nass to estimate the distribution for the null hypothesis 
of multinomial backtests.

Under the assumption that \( L \overset{d}{\equiv} M \), Kratz & al. (2018) show that
\((O_0, O_1, \ldots, O_{m+1})\)
possesses a multinomial distribution with \( n \) trials

This is used as the null hypothesis for the backtest
Multinomial tests are reviewed in detail in Cai & Krishnamoorthy (2006) and compared numerically:

Very famous is Pearson's \(\chi^2\)-test published in 1900
The cumulated relative squared deviations of the observations from the mean

\(
S_{m+1} := \sum_{k=0}^{m+1} \frac{(O_k - np_k)^2}{np_k},
\)

are for large \( n \) approximately \(\chi^2_{m+1}\)-distributed
The asymptotic test rejects the null hypothesis when \( S_{m+1} \) is too large

Nass (1959) suggests a finite sample correction of Pearson's \(\chi^2\)-test
Instead of approximating the distribution of \(S_{m+1}\) with a \(\chi^2_{m+1}\)-distribution, he proposes to use a distribution that depends on \(n\), namely the distribution of \(\frac{1}{c}Z\) with \(Z \sim \chi^2_{\nu}\) where

\(
c = \frac{2E[S_{m+1}]}{\text{Var}(S_{m+1})}, \quad \nu = cE[S_{m+1}], \quad E[S_{m+1}] = m+1,
\)

\(
\text{Var}(S_{m+1}) = 2(m+1) - \frac{m^2 + 6m + 6}{n} + \sum_{k=0}^{m+1} \frac{1}{n \cdot p_k}
\)

Since \(c \to 1\) and \(\nu \to m+1\) as \(n \to \infty\), the distributions used in the finite sample approximation are asymptotically equal to the asymptotic distribution 
known from Pearson’s \(\chi^2\)-test
But in Nass' \(\chi^2\)-test the distribution of the approximating \(\frac{1}{c}Z\) matches for each \(n\) the first two moments of the distribution 
of \(S_{m+1}\), while \(\chi^2_{m+1}\) matches only the first moment


How can the AV@R-approach be generalized to arbitrary DRMs for multinomial backtests?

If \(g\) is left-continuous, the DRM can be expressed as
\(
\rho_g^{t-1}(M_t) = \int_{[0,1]} q_{M_t}^{t-1}(u) d\bar{g}(u), \quad \bar{g}(u) = 1 - g(1-u)
\)
for a right-continuous distribution function \(\bar{g}\) on the interval \([0, 1]\).

Let \(\bar{G} \sim \bar{g}\) be an independent random variable, \(G = 1 - \bar{G}\), then conditionally on \(M_{t-1} = L_{t-1}, \ldots, M_1 = L_1\):
\(
\rho_g^{t-1}(M_t) = \mathbb{E} \left[ q_{M_t}^{t-1}(\bar{G}) \right] = \mathbb{E} \left[ q_{M_t}^{t-1}(1 - G) \right]
\)
Exception indicators are now obtained by sampling quantile levels from \(\bar{G} = 1 - G\).
Let \(0 = \alpha_0 < \alpha_1 < \cdots < \alpha_m < \alpha_{m+1} = 1\) be a partition of \( [0,1] \) with \( g(\alpha_j) - g(\alpha_{j-1}) \neq 0 \) 
for all \( j = 1, \ldots, m+1 \).
For \( t, j \) we let \( G_{t,j} \) be independent random variables with distribution \( \mathcal{L}(G \mid G \in [\alpha_{j-1}, \alpha_j)) \)
The exception indicators that can be used for backtesting are
\[
\mathbf{1}_{t,j} = 
\begin{cases} 
1 & \text{if } L_t > q_{M_t}^{t-1}(1 - G_{t,j}) \\
0 & \text{else} 
\end{cases}
, \quad j = 1, \ldots, m+1
\]

Theorem:
If M and L possess the same law, then
\(
\mathbb{E}[\mathbf{1}_{t,j}] = \frac{\mathbb{E}\left[G \mathbf{1}_{\{G \in [\alpha_{j-1}, \alpha_j)\}}\right]}{g(\alpha_j) - g(\alpha_{j-1})}
\)

for all \( t = 1, \ldots, n \) and the random vectors \(\mathbf{1}_{t,i})_{i=1,...,m+1}, t = 1, \ldots, n\), are independent.

The numbers of breached randomized levels are \( X_t := \sum_{j=1}^{m+1} \mathbf{1}_{t,j} \), \( t \in \{1, \ldots, n\} \)
The observed cell counts
\[
O_k := \sum_{t=1}^n \mathbf{1}_{\{X_t = k\}}, \quad k = 0, 1, \ldots, m+1
\]

are the key statistics of the multinomial backtest.

Describe explicitely how to run a backtest for DRMs. Explain the power impact of
stratification of levels and discretization.

Null hypothesis \( H_0 \)
Components of \( L = (L_t)_{t=1,2,\ldots,n} \) are independent random variables with standard normal distribution \( \mathcal{N} \)
This is also the basis for the risk computation within the model

Alternative \( H_1 \)
Alternative distributions, but always scaled and shifted such that they possess expectation 0 and unit variance as in the standard normal model, i.e.,

    \[
    L_t = (\tilde{L}_t - E[\tilde{L}_t]) / \sqrt{\text{Var}(\tilde{L}_t)}, \quad t = 1, 2, \ldots, n
    \]

T3: student-t with three degrees of freedom
T5: student-t with five degrees of freedom
ST: skewed-t-distribution of Fernández and Steel (1998)

Size and power of the asymptotic tests are obtained by simulating the stochastic process for the null hypothesis and the alternatives. 
This is repeated \( N \) times to obtain the estimator
\[
\frac{1}{N} \sum_{i=1}^{N} \mathbf{1}_{\{H_0 \text{ is rejected}\}} (L_{1,i}, \ldots, L_{n,i})
\]

A larger number of samples n leads to improved power
Power in our method does not monotonously increase in the stratification parameter m
In particular, too many strata, i.e., large m, is not the best choice - randomization of levels is more
successful than discritization!
